\section{Eight rank-nine exclusions from central restrictions}
All fusion categories in this section are over $\mathbb C$. The independent
signed-sphericalization certificates show that a categorification of any of
the eight rings considered here has a Galois conjugate with a spherical
structure whose simple dimensions are its Frobenius--Perron dimensions.
It therefore suffices to exclude pseudo-unitary categorifications.

Write $F\colon\mathcal Z(\mathcal C)\to\mathcal C$ for the forgetful
functor and $I$ for its adjoint. Put $D=\operatorname{FPdim}(\mathcal C)$.
For each irreducible representation $E$ of $K_0(\mathcal C)\otimes\mathbb C$,
Ostrik's Theorem~2.13 in \emph{Pivotal fusion categories of rank~3}
(Moscow Math. J. \textbf{15} (2015), 373--396;
\texttt{arXiv:1309.4822}) gives a simple object $A_E$ of the center such that
\[
 \dim A_E=D/f_E,\qquad [F(A_E):\mathbf1]=\dim E.
\]
Here $f_E$ is the formal codegree. The class $[F(A_E)]$ is central in
$K_0(\mathcal C)$, since a half-braiding gives isomorphisms
$F(A_E)\otimes X\cong X\otimes F(A_E)$.
For one-dimensional characters $\chi,\psi$, the modular character formula
and symmetry of the unnormalized $S$-matrix also give
\begin{equation}\label{eq:central-reciprocity}
 \frac D{f_\chi}\,\chi([F(A_\psi)])
 =\frac D{f_\psi}\,\psi([F(A_\chi)]).
\end{equation}
Indeed the character of $K_0(\mathcal Z(\mathcal C))$ obtained by composing
$\chi$ with $F$ is $[B]\mapsto S_{A_\chi,B}/\dim A_\chi$.
This is the correspondence preceding Ostrik's Theorem~2.13.

\subsection{A dimension-lattice obstruction}
Let $L$ be the additive subgroup of $\mathbb R$ generated by the simple
Frobenius--Perron dimensions of $\mathcal C$. Every $D/f_E$ belongs to $L$:
it is the dimension of the object $F(A_E)$. This gives a necessary
condition stronger than algebraic integrality of $D/f_E$.

For each of
\[
 R9\_1\_15,\quad R9\_1\_68,\quad R9\_1\_88,\quad R9\_1\_116,
\]
the dimension multiset is $(1^4,2,a^2,b^2)$, with
$a\simeq3.603875$ and $b\simeq4.493959$; their minimal polynomials are
\[
 a^3-2a^2-8a+8=0,\qquad b^3-4b^2-4b+8=0.
\]
The three numbers $1,a,b$ are rationally independent, and
$D=24+4a+8b$. The formal codegree $8$ occurs. Hence a hypothetical
categorification would satisfy
\[
 3+\tfrac12a+b=\frac D8\in\mathbb Z+\mathbb Za+\mathbb Zb,
\]
contradicting rational independence. The verifier reconstructs the
Frobenius--Perron character, the degree-three field, the independence,
and the codegree directly from each input tensor.

\subsection{A three-character obstruction}
Put $d=2+\sqrt6$, so $d^2=4d+2$. Each of
\[
 R9\_1\_17,\quad R9\_1\_83,\quad R9\_1\_122
\]
has four invertibles, one simple $Y$ of dimension $2$, and four simples
of dimension $d$. Its global dimension is $D=16(1+d)$.
There are two one-dimensional characters $\chi_+,\chi_-$ of codegree
$16$ and one character $\psi$ of codegree $8$ with the following
properties. Both $\chi_\pm$ take value $1$ on every invertible and $-2$
on $Y$. The four dimension-$d$ simples split into two pairs, on which
$\chi_+$ takes values $\sqrt2$ and $-\sqrt2$, respectively;
$\chi_-=-\chi_+$ on these four objects. The character $\psi$ vanishes
on one pair and is nonzero on both members of the other pair.
These assertions are verified as exact character identities against
all fusion products; no inferred character-table completeness is used.

The central objects $A_+,A_-$ corresponding to $\chi_+,\chi_-$ have
dimension $1+d$ and contain the unit with multiplicity one under $F$.
Rational independence of $1,d$ therefore forces
\[
 [F(A_+)]=1+X,\qquad [F(A_-)]=1+Z
\]
for dimension-$d$ simples $X,Z$. Equation~\eqref{eq:central-reciprocity}
for this pair gives $\chi_+(Z)=-\chi_+(X)$, so $X,Z$ lie in opposite
pairs. The central object $A_\psi$ has dimension $2+2d$ and unit
multiplicity one. Thus its restriction has exactly two invertible
summands and two dimension-$d$ summands, counting multiplicities,
and no copy of $Y$. Consequently
\[
 (\chi_++\chi_-)([F(A_\psi)])=4.
\]
Summing \eqref{eq:central-reciprocity} for $(\chi_+,\psi)$ and
$(\chi_-,\psi)$, and using $D/8=2D/16$, gives
\[
 4=2\bigl(2+\psi(X)+\psi(Z)\bigr).
\]
Hence $\psi(X)+\psi(Z)=0$. But exactly one of these two values is
nonzero, since $X,Z$ lie in opposite pairs. This contradiction excludes
all three rings. It uses only three central simple objects and three
characters, with no enumeration of induction matrices.

\subsection{A central-support obstruction}
The remaining ring $R9\_1\_120$ has the same dimension multiset
$(1^4,2,d^4)$ and global dimension $16(1+d)$, but is noncommutative.
It has a one-dimensional character of formal codegree $16$ (its values
on the dimension-$d$ simples are $\pm i\sqrt2$).
The corresponding central simple again restricts to $1+X$ for a single
dimension-$d$ simple $X$. None of the four possible $X$ is central in
the fusion ring. Indeed, in the original zero-based input order,
$X_1,X_3$ fail to commute with $X_2$, while $X_2,X_4$ fail to commute
with $X_1$. Thus $1+X$ cannot be the restriction of a central object.
The certificate records a differing coefficient for each of the four
commutators. This proves the fourth quadratic-family exclusion.

\paragraph{Verification files.}
\nolinkurl{cubic_lattice.py} verifies the four lattice obstructions;
\nolinkurl{quadratic_exact.py} verifies the three-character and
central-support obstructions. Both reconstruct their input tensors
from the supplied archives and bind their receipts to SHA-256 hashes.
Neither final verifier uses floating-point decisions, optimization,
or the exploratory induction-matrix computations.
